\chapter{Conventions and Notation}
\label{chap:conventions}

The following chapters use several sign conventions and two distinct ways in
which a map can act.  This short chapter fixes those choices once, so that later
formulae can be read without silently changing conventions.  The choices are
those used throughout this document and are compatible with the accelerator
Lie-map framework of Refs.~\cite{DragtFinn1976,DragtForest1983,Forest_book1,Forest_book2}.

\section{Vectors, matrices, and canonical coordinates}

Scalars are written in ordinary italic mathematics.  Phase-space vectors are
written in bold, for example $\vec z$, and matrices are denoted by capital
letters.  The superscripts $T$ and $\dagger$ denote transpose and Hermitian
conjugation, respectively; $\ii$ denotes the imaginary unit.  Calligraphic
symbols such as $\mathcal M$ denote maps when this helps distinguish them
from their Jacobian matrices.  The differential of a map at $\vec z$ is written
$D\mathcal M_{\vec z}$; in a coordinate chart it is represented by the usual
Jacobian matrix.  Repeated indices are summed only where this is stated or where
the Einstein convention is explicitly being used; most sums in the accelerator
calculations are written out.

Canonical coordinates are ordered in conjugate pairs,
\begin{equation}\label{eq:conv-coordinates}
    \vec z = (q_1,p_1,\ldots,q_n,p_n)^T\,.
\end{equation}
With this ordering, define
\begin{equation}\label{eq:conv-S}
    S = \diag(S_2,\ldots,S_2)\,,
    \qquad
    S_2 = \begin{pmatrix}
        0 & 1\\
        -1 & 0
    \end{pmatrix}\,.
\end{equation}
In canonical coordinates the symplectic form is therefore
\begin{equation}\label{eq:conv-symplectic-form}
    \omega = \sum_{j = 1}^n \dee q_j\wedge\dee p_j\,,
    \qquad
    \omega(\vec u,\vec v) = \vec u^T S\vec v\,.
\end{equation}
A linear map $\vec z\mapsto R\vec z$ is symplectic when
$R^T S R = S$; for a nonlinear canonical map the same condition holds for its
Jacobian at every phase-space point.

\section{Poisson bracket, Hamiltonian flow, and Lie maps}

The Poisson bracket is defined with the first argument differentiated with
respect to $q_j$ in the first term,
\begin{equation}\label{eq:conv-poisson}
    \poissonfull{f}{g} = \sum_{j = 1}^n\left(
        \partder{f}{q_j}\partder{g}{p_j}
        - \partder{f}{p_j}\partder{g}{q_j}
    \right)\,.
\end{equation}
Thus $\poisson{q_i}{p_j} = \delta_{ij}$.  The Hamiltonian vector field is
defined by
\begin{equation}\label{eq:conv-hamiltonian-vector-field}
    \iota_{X_f}\omega = \dee f\,.
\end{equation}
With the symplectic-form convention above, the bracket can equivalently be
written
\begin{equation}\label{eq:conv-poisson-geometric}
    \poisson{f}{g} = \omega(X_f,X_g) = X_g f = -X_f g\,.
\end{equation}
Thus
\begin{equation}\label{eq:conv-hamilton-equations}
    \dot q_j = \partder{H}{p_j}\,,
    \qquad
    \dot p_j = -\partder{H}{q_j}\,,
    \qquad
    \dot g = \poisson{g}{H}\,.
\end{equation}
The Lie operator places its generator in the \emph{first} Poisson-bracket
slot,
\begin{equation}\label{eq:conv-lie-operator}
    \Lie{f}\,g \equiv \poisson{f}{g}\,.
\end{equation}
Consequently $X_f g = \poisson{g}{f} = -\Lie{f}g$.  An expression such as
$\e^{\lie{f}}$ always means the exponential of the \emph{operator} $\Lie{f}$;
it is not the scalar function $\e^f$.  With these definitions, the forward
Hamiltonian point flow is
\begin{equation}\label{eq:conv-forward-flow}
    \Phi_H^t(\vec z) = \e^{-t\,\lie{H}}\vec z\,.
\end{equation}
This minus sign is intentional.  It follows from the ordering in
Eq.~\eqref{eq:conv-lie-operator}, rather than from a different sign in
Hamilton's equations.  In particular, when a Hamiltonian point map is denoted
by $\mathcal E_H$, this document uses
\begin{equation}\label{eq:conv-EH}
    \mathcal E_H(\vec z) = \e^{-\lie{H}}\vec z\,.
\end{equation}
Accelerator Lie-map references often write the same finite transformation as
$\e^{\lie{f}}$ with a \emph{map generator} $f$.  There is no sign conflict:
for an autonomous physical Hamiltonian propagated for a parameter interval
$L$, one simply has
\begin{equation}\label{eq:conv-map-generator}
    \e^{-L\,\lie{H}} = \e^{\lie{f}}\,,
    \qquad f = -L H\,.
\end{equation}
Keeping the physical Hamiltonian $H$ and the Lie-map generator $f$ distinct
is useful when comparing conventions across the literature.

\section{Point maps and substitution operators}

A point map $\mathcal A:M\to M$ acts on phase-space points.  Its associated
substitution operator acts on functions by pullback,
\begin{equation}\label{eq:conv-substitution}
    \widehat{\mathcal A}\,g = g\circ\mathcal A\,.
\end{equation}
If a particle encounters $\mathcal A$ first and $\mathcal B$ second, then
\begin{equation}\label{eq:conv-map-order}
    \vec z_{\rm out} = (\mathcal B\circ\mathcal A)(\vec z_{\rm in})\,,
    \qquad
    \widehat{\mathcal B\circ\mathcal A} = \widehat{\mathcal A}\,\widehat{\mathcal B}\,.
\end{equation}
Thus point-map composition and substitution-operator multiplication appear in
opposite orders.  Whenever the distinction matters, this document keeps the
composition symbol or the map argument explicit.

\section{Normal-form phase convention}

For each linearly normalised degree of freedom, the action--angle variables are
chosen as
\begin{equation}\label{eq:conv-action-angle}
    J_j = \frac{Q_j^2 + P_j^2}{2}\,,
    \qquad
    Q_j = \sqrt{2J_j}\cos\phi_j\,,
    \qquad
    P_j = -\sqrt{2J_j}\sin\phi_j\,.
\end{equation}
The corresponding complex phasor is
\begin{equation}\label{eq:conv-phasor}
    a_j = \frac{Q_j-\ii P_j}{\sqrt 2} = \sqrt{J_j}\,\e^{\ii\phi_j}\,.
\end{equation}
A positive phase advance $\mu_j$ therefore acts as
\begin{equation}\label{eq:conv-phase-advance}
    a_j \longmapsto \e^{\ii\mu_j}a_j\,,
    \qquad
    \phi_j \longmapsto \phi_j + \mu_j\,,
    \qquad
    \nu_j = \frac{\mu_j}{2\pi}\,.
\end{equation}
A one-turn eigenvalue determines $\mu_j$ modulo $2\pi$ and the tune modulo an
integer.  Continuous phase or tune branches must therefore be chosen when
parameters or observation locations are varied.

\section{Series order and external parameters}

A dot denotes differentiation with respect to the continuous flow parameter in
context (physical time, longitudinal position, or an auxiliary Lie-flow parameter),
whereas an integer subscript $n$ denotes the discrete turn number.

The symbol $\mathcal O(\|\vec z\|^r)$ denotes terms of coordinate degree at
least $r$ about the stated expansion point.  A homogeneous Hamiltonian of
degree $r$ generates a first coordinate change of degree $r-1$, since one
Poisson bracket differentiates the generator once.  Normal-form calculations
are understood locally as formal power-series calculations, or as finite
truncations of such series, unless convergence is stated explicitly.

Quantities such as a magnet strength or a prescribed momentum offset may be
carried as external Taylor parameters.  Unless a parameter is explicitly
promoted to a canonical dynamical coordinate with its own conjugate momentum,
it is held fixed by the Poisson bracket.  Complex phasors are a computational
basis; on real physical phase space their conjugates satisfy
$a_j^* = \overline{a_j}$.
